% 01-resumen.tex — Executive Summary (2 pages)
\chapter{Executive Summary}
\label{ch:summary}

\section{What is Doki?}
\label{sec:what-is-doki}

Doki is an open-source container ecosystem designed for Android and
Linux environments. Unlike traditional container runtimes that require
root access and specialized kernel features, Doki operates entirely in
user space using \texttt{proot} for syscall interception and filesystem
virtualization. The system consists of four binaries---\texttt{doki}
(CLI), \texttt{dokid} (daemon), \texttt{doki-compose} (orchestration),
and \texttt{doki-init} (initialization)---that together provide a
complete container lifecycle management solution in approximately
28.3\,MB of statically linked Go binaries.

\section{The Problem}
\label{sec:problem}

Android presents unique challenges for containerization. The Linux
kernel underlying Android restricts access to critical system calls,
the \texttt{iptables} subsystem is typically unavailable without root,
and SELinux policies block common networking operations such as binding
to port 53. Traditional container runtimes like Docker and containerd
are designed for server environments and cannot operate under these
constraints.

The mobile development community lacks a tool that provides Docker-like
functionality on Android devices. Developers working with containers
on phones and tablets must either use simplified chroot environments
or rely on cloud-based solutions that introduce latency and connectivity
dependencies.

\section{Key Achievements}
\label{sec:achievements}

Doki v0.9.3 delivers the following capabilities:

\begin{itemize}[leftmargin=*, itemsep=4pt]
  \item \textbf{12 isolation levels} --- from WASM sandboxes to
        hardware-level microVMs, with automatic detection of the
        strongest mode available on the host.
  \item \textbf{244 CLI commands} --- covering container management,
        image operations, networking, orchestration, and mesh
        connectivity.
  \item \textbf{DokiLink-Lite mesh networking} --- peer-to-peer
        encrypted communication using TLS 1.3 with optional NaCl
        secretbox payload encryption.
  \item \textbf{Rootless operation} --- full functionality without
        root access, using \texttt{proot} for syscall interception
        and \texttt{fuse-overlayfs} for filesystem virtualization.
  \item \textbf{Docker Compose compatibility} --- support for
        Dockerfile syntax, Docker Compose files, and BuildKit-style
        multi-stage builds.
  \item \textbf{Multi-architecture support} --- ARM64 and ARMv7
        binaries for Android, Linux, and macOS (Darwin).
\end{itemize}

\section{Impact}
\label{sec:impact}

Doki enables developers to run containerized workloads on Android
devices without root access, cloud connectivity, or specialized
hardware. This opens new possibilities for:

\begin{itemize}[leftmargin=*, itemsep=4pt]
  \item \textbf{Mobile development environments} --- running
        compilers, linters, and test suites directly on phones.
  \item \textbf{Edge computing} --- deploying containerized services
        on mobile devices at the network edge.
  \item \textbf{Education} --- teaching container concepts using
        widely available Android hardware.
  \item \textbf{Field operations} --- running isolated tools in
        environments without traditional server infrastructure.
\end{itemize}

\section{Document Structure}
\label{sec:doc-structure}

This report is organized as follows. \cref{ch:architecture} describes
the overall system architecture. \cref{ch:components} details the
individual subsystems. \cref{ch:isolation} covers the 12 isolation
levels. \cref{ch:networking} explains the networking stack and
DokiLink-Lite. \cref{ch:storage} describes image storage and layer
management. \cref{ch:build} covers the build system and Dockerfile
support. \cref{ch:metrics} presents performance benchmarks.
\cref{ch:roadmap} outlines future development plans.

\begin{figure}[htbp]
  \centering
  \begin{tikzpicture}[
    node distance=0.5cm,
    comp/.style={
      rectangle, rounded corners=3pt,
      draw=nodeblue!60, fill=fillblue,
      minimum width=2.8cm, minimum height=1.1cm,
      font=\sffamily\small, text=textdark,
      align=center, inner sep=5pt
    },
    compAlt/.style={
      comp, draw=nodeorange!60, fill=fillorange
    },
    compTeal/.style={
      comp, draw=nodelteal!60, fill=fillteal
    },
    arr/.style={
      -{Stealth[length=3mm, width=2mm]},
      draw=nodeblue!60, line width=0.7pt
    }
  ]
    % Top row
    \node[comp] (cli) {doki\\{\scriptsize\color{textgray} CLI Client}};
    \node[comp, right=2cm of cli] (daemon) {dokid\\{\scriptsize\color{textgray} Daemon}};
    \node[compAlt, right=2cm of daemon] (compose) {doki-compose\\{\scriptsize\color{textgray} Orchestration}};

    % Bottom row
    \node[compTeal, below=1cm of daemon] (runtime) {Runtime\\{\scriptsize\color{textgray} 12 Isolation Levels}};
    \node[comp, below=1cm of runtime] (kernel) {Linux Kernel\\{\scriptsize\color{textgray} / proot}};

    % Arrows
    \draw[arr] (cli) -- (daemon);
    \draw[arr] (daemon) -- (compose);
    \draw[arr] (daemon) -- (runtime);
    \draw[arr] (runtime) -- (kernel);

    % Region labels
    \node[annot, above=0.2cm of cli.north west] {\textbf{User Space}};
    \node[annot, above=0.2cm of daemon.north] {\textbf{Daemon}};
    \node[annot, above=0.2cm of runtime.north] {\textbf{Runtime}};
  \end{tikzpicture}
  \caption{High-level overview of the Doki ecosystem components.}
  \label{fig:overview}
\end{figure}

\clearpage
